%  ::::::::::::::::::::::::::::::::::::::::::::::::::::::::::::::::::::::::::::::::::::::::::::::
%  ::::::::::::::::::::::::::::::::::::::::::::::::::::::::::::::::::::::::::::::::::::::::::::::
%  ::::::::::::::::::::::::::::::::::::::::::::::::::::::::::::::::::::::::::::::::::::::::::::::
%  ::::::::::::::::::::::::::::::::::::::::::::::::::::::::::::::::::::::::::::::::::::::::::::::
%  ::::::::::::::::::::::::::::::::::::::::::::::::::::::::::::::::::::::::::::::::::::::::::::::
%  ::::::::::::::::::::::::::::::::::::::::::::::::::::::::::::::::::::::::::::::::::::::::::::::
%  ::::::::::::::::::::::::::::::::::::::::::::::::::::::::::::::::::::::::::::::::::::::::::::::
%  ::::::::::::::::::::::::::::::::::::::::::::::::::::::::::::::::::::::::::::::::::::::::::::::
%  ::::::::::::::::::::::::::::::::::::::::::::::::::::::::::::::::::::::::::::::::::::::::::::::
%  ::::::::::::::::::::::::::::::::::::::::::::::::::::::::::::::::::::::::::::::::::::::::::::::
%  ::::::::::::::::::::::::::::::::::::::.................:::::::::::::::::::::::::::::::::::::::
%  ::::::::::::::::::::::::::::::::.............................::::::::::::::::::::::::::::::::
%  ::::::::::::::::::::::::::::......................................:::::::::::::::::::::::::::
%  ::::::::::::::::::::::::......................*%:....................::::::::::::::::::::::::
%  ::::::::::::::::::::::.......................+@@@-......................::::::::::::::::::::::
%  ::::::::::::::::::::........................+@@@@@:.......................:::::::::::::::::::
%  ::::::::::::::::::.........................=@@@@@@@:........................:::::::::::::::::
%  ::::::::::::::::..........................:@@@@@@@@@-........................:::::::::::::::
%  :::::::::::::::..........................-@@@@@@@@@@@=.........................:::::::::::::
%  :::::::::::::...........................=@@@@@@@@@@@@@-.........................::::::::::::::
%  ::::::::::::...........................-@@@@@@@@@@@@@@@..........................:::::::::::
%  :::::::::::............................:%@@@@@@@@@@@@@+...........................:::::::::
%  ::::::::::..............................=@@@@@@@@@@@@%:............................:::::::::
%  ::::::::::...............................*@@@@@@@@@@@=..............................::::::::
%  :::::::::................................:@@@@@@@@@@%:...............................::::::
%  ::::::::..................................*@@@@@@@@@-................................::::::::
%  ::::::::..................:@@+:...........:@@@@@@@@@.............:+-..................:::::::
%  :::::::...................*@@@@@@*-:.......%@@@@@@@+........:-*@@@@@..................:::::::
%  :::::::..................:@@@@@@@@@@@%:....*@@@@@@@:....:=%@@@@@@@@@=.................:::::::
%  :::::::..................*@@@@@@@@@@@@#....=@@@@@@@....:*@@@@@@@@@@@#..................::::::
%  :::::::.................:@@@@@@@@@@@@@@-...=@@@@@@@....*@@@@@@@@@@@@@:.................::::::
%  :::::::.................*@@@@@@@@@@@@@@@:..=@@@@@@#...+@@@@@@@@@@@@@@=.................::::::
%  :::::::................:@@@@@@@@@@@@@@@@*..=@@@@@@#..+@@@@@@@@@@@@@@@+.................::::::
%  :::::::................=@@@@@@@@@@@@@@@@@-.#@@@@@@@.-@@@@@@@@@@@@@@@@*................:::::::
%  :::::::...............:#@@@@@@@@@@@@@@@@@*.@@@@@@@@:@@@@@@@@@@@@@@@@@%:...............:::::::
%  ::::::::..............:*@@@@@@@@@@@@@@@@@@@@@@@@@@@@@@@@@@@@@@@@@@@@@%:...............:::::::
%  ::::::::................:*@@@@@@@@@@@@@@@@@@@@@@@@@@@@@@@@@@@@@@@@@@@-...............::::::::
%  :::::::::.................:=#@@@@@@@@@@@@@@@@@@@@@@@@@@@@@@@@@@@@@%-.................::::::::
%  ::::::::::....................:#@@@@@@@@@@@@@@@@@@@@@@@@@@@@@@@@=...................::::::::::
%  ::::::::::.......................:*@@@@@@@@@@@@@@@@@@@@@@@@@#-.....................:::::::::
%  :::::::::::.........................:=@@@@@@@@@@@@@@@@@@*:........................:::::::::::
%  ::::::::::::......................:=%@@@@@@@@@@@@@@@@@@@@#:......................::::::::::::
%  :::::::::::::.............+#%@@@@@@@@@@@@@@%-::*-.:%@@@@@@@@%=:.................::::::::::::::
%  :::::::::::::::...........:#@@@@@@@@@@@#--+%@@@@@@@#=:=%@@@@@@@@@@-............::::::::::::::::
%  ::::::::::::::::............-@@@@@@+-=#@@@@@@@@@@@@@@@@#=-=#@@@@*:............::::::::::::::::
%  ::::::::::::::::::...........:==:...-@@@@@@@@@@@@@@@@@@@@:...:=-............:::::::::::::::::
%  :::::::::::::::::::...................@@@@@@@@@@@@@@@@@-..................::::::::::::::::::::
%  ::::::::::::::::::::::................:#@@@@@@@@@@@@@*:.................::::::::::::::::::::::
%  ::::::::::::::::::::::::...............:*@@%+-.:=#@%-................::::::::::::::::::::::::
%  ::::::::::::::::::::::::::::.............:........................:::::::::::::::::::::::::::
%  :::::::::::::::::::::::::::::::...............................:::::::::::::::::::::::::::::::
%  ::::::::::::::::::::::::::::::::::::.....................:::::::::::::::::::::::::::::::::::
%  ::::::::::::::::::::::::::::::::::::::::::::::::::::::::::::::::::::::::::::::::::::::::::::::
%  ::::::::::::::::::::::::::::::::::::::::::::::::::::::::::::::::::::::::::::::::::::::::::::::
%  ::::::::::::::::::::::::::::::::::::::::::::::::::::::::::::::::::::::::::::::::::::::::::::::
%  ::::::::::::::::::::::::::::::::::::::::::::::::::::::::::::::::::::::::::::::::::::::::::::::

% HyperTensor Paper XVIII: The Bridge Protocol
\documentclass[11pt,a4paper]{article}
\usepackage[margin=1in]{geometry}
\usepackage{amsmath,amssymb,amsthm}
\usepackage{hyperref,graphicx,booktabs,longtable,enumitem}
\usepackage{setspace}
\onehalfspacing
\hypersetup{colorlinks=true,linkcolor=blue,urlcolor=blue,citecolor=blue}

\title{The Bridge: A Unified Protocol for the Riemann Hypothesis\\\large HyperTensor Paper XVIII}
\author{William Ken Ohara Stewart\\NagusameCS Independent Research\\\texttt{https://github.com/NagusameCS/HyperTensor}}
\date{May 5, 2026}

\pagestyle{empty}

\begin{document}
\maketitle

\begin{abstract}
We compose Papers XVI (AGT) and XVII (ACM) into a unified 5-step protocol
that detects, verifies, and validates zeros of the Riemann zeta function
on the critical line. The protocol chain is: AGT detection via the
difference operator $D(s)$, ACM verification of $Z_2$ invariance, Safe
OGD exploration constrained to the critical-line neighbourhood, TEH
exclusion of off-critical candidates, and contradiction closure. The
protocol correctly identifies all 105 known non-trivial zeros of $\zeta(s)$,
which lie on $\mathrm{Re}(s)=1/2$ by classical results (Hardy 1914, Selberg,
van de Lune et al., Platt et al.).

\textbf{What this is not:} The construction is definitional---$D_0(s)=2\sigma-1$ is zero iff $\sigma=1/2$ by design. The ``remaining step'' of proving $\zeta(s)=0 \implies D(s)=0$ is the Riemann Hypothesis itself, restated in different notation. Substituting $D$'s definition, this is ``prove every nontrivial zero has $\sigma=1/2$,'' which is RH. This paper does NOT reduce RH to an attainable next-step theorem. The $J\approx 1-10^{-315}$ figure treats the 105 known zeros as independent jurors, but they are consequences of the same classical theorem-and-computation pipeline. This is a geometric visualisation of the functional equation's $Z_2$ symmetry, not a proof-search protocol.
\end{abstract}



% =================================================================
\section{Introduction: A Geometric Visualisation of $Z_2$ Symmetry}
% =================================================================

The Riemann Hypothesis has resisted proof for 165 years. Papers XVI and XVII of the HyperTensor framework developed two geometric approaches to visualising the functional equation's $Z_2$ symmetry: AGT uses a hand-designed prime-based feature map, and ACM learns the involution in a neural latent space. Paper XVIII composes them into a unified visualisation pipeline.

\textbf{Important caveat:} The construction is definitional. Feature coordinate $f_0(s)=\sigma$ (the real part of $s$) is explicitly encoded, so $D_0(s)=2\sigma-1$ is zero iff $\sigma=1/2$ by design. The 100\% detection rate confirms this definition; it does not constitute evidence for the Riemann Hypothesis. The ``remaining step''---proving $\zeta(s)=0 \implies D(s)=0$ for all zeros---is the Riemann Hypothesis itself, restated. Substituting $D$'s definition, this step reads ``prove every nontrivial zero has $\sigma=1/2$,'' which is RH. The framework renames the problem; it does not reduce it.

What follows is a geometric visualisation of the functional equation's $Z_2$ symmetry, with no claim toward proof.

% =================================================================
\section{The 5-Step Pipeline}
% =================================================================

\subsection{Step 1: AGT Detect}

Scan candidate $s$-values in the critical strip $0 < \mathrm{Re}(s) < 1$. For each candidate, compute the difference operator $D(s) = f(s) - f(\iota(s))$ using the AGT feature map. Candidates with $D(s) \approx 0$ are candidate zeros on the critical line. Candidates with $D(s) > 0$ are off-critical and rejected.

At 50,000 primes: 100\% detection of critical-line points, 0\% false positives.

\subsection{Step 2: ACM Verify}

Pass candidate zeros through the ACM embedder. Verify that $h(\iota(s)) \approx h(s)$ (fixed-point property) and that $\iota^2(s) \approx s$ (involution property). Candidates satisfying both are geometrically confirmed on the critical line.

At 105 zeros: fixed-point error 0.008, involution error 0.009.

\subsection{Step 3: Safe OGD Explore}

Use the Safe OGD projector $P_{\text{safe}} = I - Q_f Q_f^T$ from Paper XIII to search the neighborhood of each candidate zero. The projector constrains the search to the critical line, preventing drift into off-critical territory. This verifies that the candidate zero lies in a connected component of critical-line points.

\subsection{Step 4: TEH Exclude}

Apply the Tangent Eigenvalue Harmonics detector. Any point with $\mathrm{TEH}(h(s)) > 0$ has forbidden-subspace activation and is geometrically excluded from being a true zero. The TEH threshold is calibrated per model but typically $\sim$12\%.

At 105 zeros: 0 false positives (0/20 critical-line points flagged).

\subsection{Step 5: Contradiction Closes}

What the protocol can do: If a point $z$ satisfies $\zeta(z) = 0$ but fails the geometric checks (AGT $D(z) \neq 0$ or ACM $\mathrm{TEH}(z) > 0$), we have identified a counterexample to the hypothesis that ALL zeros lie on the critical line. The geometric framework would correctly detect this counterexample because the feature map is faithful by construction: $D(s) \neq 0$ for any $s$ with $\mathrm{Re}(s) \neq 0.5$.

What the protocol cannot do (the remaining gap): The converse direction — proving that every true zero MUST satisfy $\mathrm{Re}(z) = 0.5$ — requires establishing that the feature map's encoding of prime-based information is a complete representation of the zeta function's zero structure. This is the analytic bridge: proving $\zeta(s)=0 \implies D(s)=0$ for all zeros via the von Mangoldt explicit formula. The AGT feature map is faithful by construction (Step 1 works because $\sigma$ is explicitly encoded), but the step from ``the explicit formula connects zeros to primes'' to ``therefore $D(s)=0$ for all zeros'' requires rigorous analytic number theory. This is the theorem specified in the Mathematician Handoff.

% =================================================================
\section{Protocol Validation}
% =================================================================

\subsection{105-Zero Test}

\begin{itemize}[nosep]
\item 105 known zeros tested: 105/105 detected on the critical line
\item 0 false negatives: all zeros correctly identified
\item 0 false positives: no off-critical points misidentified as zeros
\item Jury confidence: $J = 1 - (1-0.999)^{105} \approx 1 - 10^{-315}$
\end{itemize}

\subsection{Meta-Jury Validation (3,713 Points)}

The \texttt{jury\_bridge.py} script (May 5, 2026) tested a comprehensive grid:

\begin{itemize}[nosep]
\item 713 points on critical line ($\sigma=0.5$): $D(s) = 0$ for 713/713, 100\%
\item 3,000 points off critical line ($\sigma \neq 0.5$): $D(s) > 0$ for 3,000/3,000, 100\%
\item Correlation: $\mathrm{Pearson}(D(s), |\sigma-0.5|) = 1.0000$
\item The encoding $D(s) = 0 \iff \mathrm{Re}(s) = 0.5$ is perfectly validated
\end{itemize}

% =================================================================
\section{The Explicit Formula Bridge}
% =================================================================

The von Mangoldt explicit formula is the key that connects everything:

$$\sum_{\rho} \frac{x^\rho}{\rho} = x - \sum_{n \le x} \Lambda(n) - \frac{\zeta'(0)}{\zeta(0)} - \frac{1}{2}\log(1-x^{-2})$$

where $\Lambda(n) = \log p$ if $n = p^k$ (and 0 otherwise), and the sum on the left runs over all non-trivial zeros $\rho$.

This formula explicitly connects zeta zeros (left side) to the prime number distribution (right side, via the von Mangoldt function $\Lambda$). Since the AGT feature map $f(s)$ encodes prime-number relationships, the explicit formula provides the analytic bridge: it guarantees that zeros of $\zeta(s)$ carry specific prime-related information, and that information is exactly what $f(s)$ captures.

\subsection{What the Remaining Step Actually Is}

\begin{center}
\fbox{\parbox{0.9\textwidth}{
The ``remaining step'' is: prove $\zeta(s)=0 \land 0<\mathrm{Re}(s)<1 \implies D(s)=0$.

Substituting $D(s)=f(s)-f(\iota(s))$ and noting $D_0(s)=2\sigma-1$, this is: prove $\zeta(s)=0 \land 0<\mathrm{Re}(s)<1 \implies \mathrm{Re}(s)=1/2$.

\textbf{This is the Riemann Hypothesis.} The framework renames it; it does not reduce it.
}}
\end{center}

The explicit formula connects zeta zeros to prime numbers. The AGT feature map encodes prime-number relationships. But the step from ``the explicit formula connects zeros to primes'' to ``therefore every zero must satisfy $f(s)=f(\iota(s))$'' has not been established, and establishing it would constitute a proof of RH. The computational verification (105/105 zeros detected) confirms that the 105 \textit{known} zeros---which are already known to lie on the critical line by classical results---are correctly identified. It does not provide evidence about unknown zeros.

% =================================================================
\section{Limitations}
% =================================================================

The construction is definitional and does not constitute progress toward a proof of RH:

\begin{enumerate}[nosep]
\item \textbf{Definitional detection.} $f_0(s) = \sigma$ is the first feature coordinate. $D_0(s) = 2\sigma-1$, which is zero iff $\sigma=1/2$ by definition. The 100\% detection rate on 3,713 test points confirms this definition, not a discovery.
\item \textbf{Rank-1 is forced.} Since coordinates 1--31 of $D$ are identically zero on all of $\mathbb{C}$ (not just on zeros), any non-zero values of $D$ live in the span of coordinate 0. The rank-1 observation follows from the feature design.
\item \textbf{Jury confidence is not inferential.} The 105 known zeros lie on $\mathrm{Re}(s)=1/2$ by classical results (Hardy 1914, Selberg, van de Lune et al., Platt et al.). They are consequences of the same theorem-and-computation pipeline, not 105 independent observations. Treating them as independent jurors and reporting $J\approx 1-10^{-315}$ misuses the jury formula by violating its independence axiom.
\item \textbf{The residual is RH.} Proving $\zeta(s)=0 \implies D(s)=0$ is the Riemann Hypothesis restated. Unfolding $D$'s definition, this is ``prove every nontrivial zero has $\sigma=1/2$,'' which is RH. The framework has not provided a reduction.
\end{enumerate}

The pedagogic value is in the $Z_2$-symmetry visualisation: the functional equation's involution $\iota(s)=1-s$, the Euler product connection to primes, and the geometric intuition that the critical line is the fixed-point set of $\iota$. These are genuine mathematical insights presented in an accessible visual form.

% =================================================================
\section{Implementation}
% =================================================================

\subsection{Script}

\texttt{scripts/jury\_bridge.py} --- runs the meta-jury validation on 3,713 points. Produces \texttt{benchmarks/jury\_bridge/faithfulness\_report.json}. Verified May 5, 2026 with Pearson $r = 1.0000$.

\subsection{Reproduction}

\begin{verbatim}
python scripts/jury_bridge.py
python scripts/agt_10k.py --scale 50000
python scripts/acm_prototype.py
python scripts/close_xvii_xviii_riemann.py
\end{verbatim}

All scripts use the \texttt{hyper\_optimize} module for efficient SVD and batched operations.

% =================================================================
\section{The Geometric Jury as Unifying Principle}
% =================================================================

Every decision in the HyperTensor framework is a jury vote. The Riemann attack was the FIRST application of this principle. The formula $J = 1 - \prod_i (1-c_i)$ is universal. For the 105 zeros, each zero serves as a juror, voting on whether candidate points lie on the critical line. The meta-jury of 3,713 grid points votes on whether the encoding is faithful. At every level, geometric evidence is aggregated into a single confidence score, and the confidence converges to certainty as the number of jurors grows.

\begin{thebibliography}{99}

\bibitem{riemann} B. Riemann, 1859.

\bibitem{mangoldt} H. von Mangoldt, ``Zu Riemann's Abhandlung `Ueber die Anzahl der Primzahlen unter einer gegebenen Grosse','' 1895.

\bibitem{edwards} H. M. Edwards, ``Riemann's Zeta Function,'' 1974.

\bibitem{odlyzko} A. M. Odlyzko, ``The $10^{20}$-th zero,'' 1992.

\bibitem{agt} W. K. O. Stewart, ``Algebraic Geometric Topology,'' Paper XVI, 2026.

\bibitem{acm} W. K. O. Stewart, ``Analytic Continuation Manifold,'' Paper XVII, 2026.

\bibitem{jury} W. K. O. Stewart, ``The Geometric Jury,'' 2026.

\bibitem{safe_ogd} W. K. O. Stewart, ``Safe OGD,'' Paper XIII, 2026.

\bibitem{teh} W. K. O. Stewart, ``COG + TEH,'' Paper XV, 2026.

\bibitem{handoff} W. K. O. Stewart, ``HANDOFF\_TO\_PHD.md,'' \texttt{docs/}, 2026.

\end{thebibliography}

\end{document}
